% quantifiers.tex -- drafted by mp (council member of ugjs/PLAN.md), 2 November 2026.
% Section 7: order of quantifiers, dependence and uniformity of every constant, boundary cases.
% Written from the statements of ugjs/PLAN.md Part 1 and from Sections 2-4 of this note as drafted
% on 2 November 2026 (notation.tex by cqft2, d2.tex by aqft); checked afterwards against d3.tex,
% d4.tex and consequences.tex (aqft) as they stood at 16:04 that day, whose statements agree.

This section states, for every result of the note and for every proof of the record that the
note replaces, the order of the quantifiers and, for every constant, what it depends on and
what it is uniform in. ``Uniform in $X$'' means that one value serves for all $X$ in the stated
range, the remaining data being fixed. The statements are those of
\S\S\ref{sec:d1}--\ref{sec:d4}, which agree with Part~1 of \texttt{ugjs/PLAN.md}.

\subsection{The fixed data}\label{sec:quant-fixed}

The interaction polynomial $\mathscr P(\xi)=\sum_{n=0}^{2p'}a_n\xi^n$, $a_{2p'}>0$,
$p'\ge1$, is fixed. In \S\S\ref{sec:d2}--\ref{sec:d4} and \S\ref{sec:cons} the standing
assumption $0\le\lambda/m_0^2<\varepsilon_{\mathscr P}$ of the record (\S2.6,
\cite[p.~14]{R}) is in force, and the note works in the working coordinates, in which the
pair $(m_0,\lambda)$ satisfies the hypotheses printed with Theorem~1.1.7, ``Let $m_0$, $K_2$
and $K_1$ be sufficiently large and let $\lambda$ be sufficiently small'' \cite[p.~594]{GJS};
Remark~3.3 of the record \cite[p.~16]{R} transports every conclusion back to the original
coordinates. \S\ref{sec:d1} uses neither: Lemma~\ref{lem:D1} holds for every $\lambda\ge0$ and
$m_0>0$. The constants fixed by these data are:
\begin{itemize}
\item[(F1)] $m_1>0$ and $C_7$, the constants of Theorem~1.1.7: ``There are constants $0<m_1$
and $O(1)$ independent of $h$, $Q_1$ and $Q_2$'' \cite[p.~594]{GJS}; the record reads
``independent of $h$'' over the class $\CE$ (Assumption~A8, Scope (d), \cite[p.~13]{R}).
They are fixed once $(m_0,\lambda,\mathscr P)$, $K_1$, $K_2$ and $\bar n$ are fixed by the
choices behind the printed hypotheses; the printed statement makes them independent of $h$,
$Q_1$ and $Q_2$, and says nothing more about their dependence.
\item[(F2)] $K_1$ and $K_2\ge1$, the constants of the printed norms (eq.~(1) of the record,
\cite[p.~10]{R}), and $\bar n=2p'+1$ (Assumption~A8(ii)).
\item[(F3)] The derived constants:
\begin{itemize}
\item $c_U=\lambda\sum_{n=1}^{2p'}|a_n|K_1K_2^{2n}n!$ (Lemma~\ref{aux:notation-cU}),
depending on $\lambda,a_1,\dots,a_{2p'},K_1,K_2$;
\item the stability constant $S_2(B)$ (Definition~\ref{aux:notation-S2def}), finite by
Theorem~C.5 of the record \cite[p.~41]{R}, which is proved conditional on the repairs R1--R10
of the V3 audit (Remark~\ref{aux:d3-V3}), and not evaluated. It depends on $(\lambda,\mathscr P,m_0)$ and on the box $B$
through its shape only (Proposition~\ref{aux:notation-S2}(iii), whose proof uses Lemma~C.4).
Independently of that proposition, the bound that the \emph{proof} of Theorem~C.5 produces
depends on $\Lambda$ only through $|\Lambda|$ and $\operatorname{diam}\Lambda$, monotonically
(V3 audit, items C5.1 and \S6; the statement of Theorem~C.5 does not record this); for a box
this is a bound depending on the shape only;
\item the ring constant
$\overline\Sigma_{3\times5}:=\sum_{j\ge1}(8j+12)e^{-m_1j}=4x(5-3x)(1-x)^{-2}$, $x=e^{-m_1}$
(\S\ref{sec:d4}), depending on $m_1$ only. It bounds the sum
$\Sigma_{3\times5}=\sum_{\Delta\not\subseteq B_{3\times5}}e^{-m_1d(\Delta,R)}$ of
Lemma~\ref{aux:notation-rings}(vi), which depends on $R$, and the sum of $e^{-m_1w_\Delta}$
over the charged squares outside $B_{3\times5}$ with the widths $w_\Delta$ of the axis-parallel
strips assigned in the proof of Theorem~\ref{thm:D4};
\item $V_*$, depending on $(\lambda,\mathscr P,m_0)$, finite (Lemma~\ref{aux:notation-wick}(iv));
the integrability of $C(x-y)^n$ on $[-1,1]^2$ is Lemma~\ref{aux:d1-cov}(ii),(iv).  $V_*$ is
taken for the constructed $V(h')$ of Lemma~\ref{aux:notation-wick}(iv) and needs no GAP.  It
bounds the record's $V(h')$ only through the identification of the record's $V(h)$ with that
object, which is the GAP recorded after Lemma~\ref{aux:notation-wick} (the record constructs no
Wick powers in $V(h)$) and enters at the record's Theorem~5.1, Step~1, and at Corollary~5.2
(\S\ref{sub:cons-51});
\item the factor $9^{p'}=(3^{p'})^2$ of the comparison $\|Y\|_4\le3^{p'}\|Y\|_2$ for $Y$ in the
chaoses of degree $\le2p'$ (proof of Theorem~C.5 of the record; Lemma~\ref{aux:notation-hyper};
the V3 audit reports it as Theorem~I.22 of Simon's book, which this drafter did not read),
depending on $p'$ only; it rests on the chaos structure that the V3 audit found unimported
(its M-3, repair R8), and it is applied to the constructed $Y=J_0V(h')$ of
Lemma~\ref{aux:notation-hyper}, which needs no GAP.
\end{itemize}
\end{itemize}
None of these depends on a cutoff $h$, on $T$, on $g$, on the position of a box, or on an
observable. They are constants of the working coordinates, and Lemma~\ref{lem:D2} and
Theorems~\ref{thm:D3}--\ref{thm:D4} say so (the rider of trap~10).

\subsection{The order of the quantifiers, statement by statement}\label{sec:quant-order}

\emph{Lemma~\ref{lem:D1}.} For all $h_1,h_2\in\CE$ and all $F\in\mathfrak F$ (indeed all
$F\in\bigcup_{p_0>1}L^{p_0}(\dmu)$, Remark~\ref{aux:d1-uses}(a)), the map
$\alpha\mapsto\ev{F}{h_\alpha}$ is $C^1$ on $[0,1]$ and \eqref{eq:d1-identity} holds. There is
no constant; no hypothesis on $\lambda$ or $m_0$ is used.  Status: proved, conditional on repair
R8 of the V3 audit (the linearity of $h\mapsto\Uc_h$ in Step~1), or, on the route written in
Step~1, on Lemma~C.4 of the record (repairs R1--R10); see the status paragraph after
Lemma~\ref{lem:D1}.

\emph{Lemma~\ref{lem:D2}}, in the form (a) of \S\ref{sec:d2}. For every finite family
$(F_a)_{a\in\mathfrak A}$ of localized monomials of type (1.1.5) or (1.1.6) with kernels in
$L^2$ and $n_i<\bar n$, every set $R$ of squares containing their localization squares, every
box $B\supseteq R$, every $h\in\CE$, and every assignment $\Delta\mapsto S_\Delta$,
$\Delta\in\mathcal D(h,B)$, of strips of widths $w_\Delta>0$ separating $\Delta$ from $R$,
\[
 \bigl|\ev{F}{h}-\ev{F}{h^B}\bigr|\le C_7\,c_U\Bigl(\sum_a\GJS{F_a}\Bigr)
 \sum_{\Delta\in\mathcal D(h,B)}e^{-m_1w_\Delta},\qquad F:=\sum_aF_a .
\]
$C_7$, $m_1$ and $c_U$ are those of (F1)--(F3); none depends on $h$, $B$, $R$, the family or the
strips. The right side depends on $h$ only through the finite set $\mathcal D(h,B)$ of the
squares outside $B$ charged by $h$, and on $F$ through the norm of the chosen decomposition
(1.1.9); the bound therefore holds with the infimum over decompositions. Of the cutoff the
lemma uses exactly $h\ge0$, $h\le1$ and compact support (Remark~\ref{aux:d2-uses}(1)): it is a
statement for $h\in\CE$ only, and nothing in the note provides it for cutoffs of variable sign,
for which $h_\alpha\notin\CE$ in general and neither Assumption~A3 nor Theorem~1.1.7 as
imported applies (trap~20 of the plan). The orientation of each application of
Theorem~1.1.7 is that of the strip $S_\Delta$, which is chosen before the bound is applied.
Lemma~\ref{lem:D2} uses Lemma~\ref{lem:D1}, and it is conditional on Lemma~C.4 of the record
(repairs R1--R10, which include R8 and so cover the condition of Lemma~\ref{lem:D1}).  The
condition enters in Step~3 of its proof, which decomposes $\Uc_{h-h^B}$ into the
$U_{\Delta,n}$.  There, Step~1 of the proof of Lemma~\ref{lem:D1} and
Lemma~\ref{aux:d1-moments}(a) give the convergence of the regularizations $\Uc^{(\kappa)}_w$,
$w=h-h^B$, and of their pieces $U^{(\kappa)}_{\Delta,n}$; Lemma~\ref{aux:notation-cutoffs}(iii)
splits $w$ into the squares; and Lemma~C.4 identifies the limit of $\Uc^{(\kappa)}_w$ with
$\Uc_h-\Uc_{h^B}$.  Repair R8 alone makes $h\mapsto\Uc_h$ linear but does not give that
identification as the note is written.

\emph{Theorem~\ref{thm:D3}.} Let $g\in\CH$, let $I\subset\R$ be a bounded interval of positive
length and let $s\in\R$; these determine $l_s$, $N_g$, $N_B$ and $B=B_s(g,I)$, with $|B|=3N_B$, and then
\[
 A:=e^{3\lambda|a_0|N_B}S_2(B),\qquad
 B':=2\,C_7\,c_U\,N_g\,\frac{e^{-m_1}}{1-e^{-m_1}} .
\]
For every finite linear combination $F$ of cylinder monomials (\S6 of the record) supported in
$I$, at the sharp time $s$, and every $T>0$,
$|\int F\,dq_{h_{T,g}}|\le A\|F\|_{L^2(d\varphi_0)}+B'\GJS F$. The order is: $g$ and $I$; then
$A$ and $B'$, which do not depend on $s$ (for $A$ through Proposition~\ref{aux:notation-S2}(iii),
conditional on R1--R10, as in Remark~\ref{aux:d3-quant}); then $s$, $F$ and $T$.
\begin{itemize}
\item $A$ depends on $\lambda$ and $a_0$ through Lemma~C.2 of the record \cite[p.~40]{R},
$Z(h^B)\ge e^{-\lambda|a_0||B|}$ (proved conditional on repair R8 of the V3 audit), and on
$(\lambda,\mathscr P,m_0)$ and $N_B$ through $S_2(B)$ (conditional on R1--R10), which depends
on the shape $3\times N_B$ only. It is uniform in $T>0$, in $s$, in
$F$ and in the position of $B$. It is \emph{not} uniform in $N_B$, which grows with the
diameter of the union of $\supp g$ and $I$, in particular with the distance from $I$ to
$\supp g$ (item m1 of the R3 verification).
\item $B'$ depends on $C_7$, $m_1$, $c_U$ and $N_g$. It is uniform in $T$, in $s$, in $F$
and in the position and the length of $I$; it depends on $g$ only through $N_g$.  The term
$B'\GJS F$ comes from Lemma~\ref{lem:D2}(a), so Theorem~\ref{thm:D3} is conditional on
Lemma~C.4 of the record (repairs R1--R10) through Lemma~\ref{lem:D2}.
\item Both depend on $g$; \S7 of the record fixes $g\in\CH$, so this is all it needs.
\item $\GJS F$ is the norm (1.1.9) of the decomposition of $F$ into its pieces localized in
the squares of row $l_s$ on the columns of \S\ref{sec:notation}
(Lemma~\ref{aux:d3-cyl}(v)), bounded with the column count $N^\square_I$ of
Remark~\ref{aux:notation-counts}; it does not depend on $s$ (Remark~6.3 of the record
\cite[p.~20]{R}).
\item Theorem~\ref{thm:D3} bounds $\int F\,dq_h$ for observables of bounded support, with
the constant $A\|F\|_{L^2}+B'\GJS F$; it is not a form of CE1, since $A$ depends on the
support of $F$ through $N_B$ and both constants depend on $g$ (trap~8).
\end{itemize}

\emph{The consequences of Theorem~\ref{thm:D3} in \S\S4 and 7 of the record.}
\begin{itemize}
\item Lemma~4.3(ii) \cite[p.~17]{R}: at each fixed $T>0$ and $g\in\CH$ the square moments
are finite by Assumption~A3 and H\"older's inequality with the exponents $(2,2)$; no
uniformity in $T$ or $g$ is claimed or used (finding F7 of the plan).
\item Lemma~7.1 \cite[p.~20]{R}: for every $g\in\CH$, every bounded interval $I$ of positive
length and every cylinder monomial $Q$ supported in $I$,
$\sup_{T>0}\|Q\hat\psi_T\|^2\le A\|Q^2\|_{L^2(d\varphi_0)}+B'\mathfrak n(Q)^2$, with $A,B'$ of
Theorem~\ref{thm:D3} at $(g,I,0)$; the same bound holds for $\|Q\Om_g\|^2$. The bound is
uniform in $T$ and depends on $g$ and $I$. Its right side replaces the printed
$C_8\,\mathfrak n(Q)^2$, so the \emph{statement} of Lemma~7.1 changes.
\item Lemmas~7.2--7.3 \cite[pp.~20--21]{R}: the square moments at the time $s$ and at the
times $\pm t/2$ obey the same bound, with $A$ and $B'$ independent of the time and of $T$.
\item Lemma~7.6 \cite[p.~22]{R}: for every $g\in\CH$, every real $f\in C_c^\infty(\R)$,
$f\ne0$, supported in an interval $I$, and every $n\ge1$,
$\E_{\mu_g}[X_f^{2n}]\le A_fK_f^{2n}(2n)!$ with $A_f:=A+B'K_1$ ($A,B'$ at $(g,I,0)$),
$K_f:=\max\{K',\sqrt2\,\sigma_f\}$, $K':=K_2^2(N^\square_I)^{1/2}\|f\|_2$ and
$\sigma_f^2:=\ip{f}{Cf}$. $A_f$ and $K_f$ do not depend on $n$, which is what Theorem~7.7 uses;
$A_f$ depends on $g$ and $I$; $K_f$ depends on $f$ (through $\|f\|_2$, $\sigma_f$ and the column
count of $I$), on $K_2$ and, through $C$, on $m_0$, and not on $g$. The radius $1/K_f$ of the
exponential moment comes from the route through the GJS norm; nothing in the note shows that it
is intrinsic to $\mu_g$ (trap~25). The statement of Lemma~7.6 changes ($C_8K_1$ and $K'$
become $A_f$ and $K_f$).
\item Theorem~7.7 \cite[p.~23]{R}: for given $e_1,\dots,e_k$ the polystrip of its Step~1 can
be taken of half-width $s_0:=\min_{j\le k,\ e_j\ne0}(2kK_{e_j})^{-1}$ ($K_f$ is defined for
$f\ne0$ only); if all $e_j=0$, then $F(z)=\E_{\mu_g}[\overline\zeta]$ is constant and every
$s_0>0$ serves, so $s_0:=+\infty$.  Indeed, by the Cauchy--Schwarz inequality and H\"older's
inequality with $k$ exponents $k$,
\[
 \E_{\mu_g}\Bigl[|\zeta|\,e^{\sum_j|\mathrm{Re}\,z_j||X_{e_j}|}\Bigr]
 \le\|\zeta\|_{L^2(\mu_g)}\prod_{j=1}^k
 \Bigl(\E_{\mu_g}e^{2k|\mathrm{Re}\,z_j||X_{e_j}|}\Bigr)^{1/(2k)},
\]
which is finite when $2k|\mathrm{Re}\,z_j|<1/K_{e_j}$ for every $j$ with $e_j\ne0$ (a factor
with $e_j=0$ is $1$). This $s_0$ depends on $e_1,\dots,e_k$ only, not on $g$. (Part~1 of the plan writes $s_0=1/K_f$, the radius of Lemma~7.6 for one $f$.)
\item Proposition~7.5 and Theorem~7.8 \cite[pp.~22--23]{R}: the gap $m_1$ is the constant of
Theorem~1.1.7 (through Lemma~7.4) and is uniform in $g$; $\widehat C_Q$ depends on $g$, and
only its finiteness is used. Their statements do not change.
\end{itemize}

\emph{Theorem~\ref{thm:D4}.} There is a constant $M_B$, determined by the fixed data
(F1)--(F3),
\[
 M_B^2:=e^{15\lambda|a_0|}S_2(B_{3\times5})\,9^{p'}V_*
 +C_7\,c_U\,\overline\Sigma_{3\times5}\cdot9K_1K_2^{8p'}(4p')!\,\lambda^2
 \Bigl(\sum_{n=0}^{2p'}|a_n|\Bigr)^2,
\]
such that for every $h\in\CE$, every $a\in\R$ and every real measurable $h'$ with
$\|h'\|_\infty\le1$ and $\supp h'\subseteq[a,a+2]$: $\int V(h')(0)^2\,dq_h\le M_B^2$. The order
is: $M_B$ first; then $h$, $a$ and $h'$.
\begin{itemize}
\item $M_B$ depends on $(\lambda,\mathscr P,m_0)$ through $C_7$, $m_1$ (also through
$\overline\Sigma_{3\times5}$), $K_1$, $K_2$, the unevaluated constant $S_2(B_{3\times5})$ of
Theorem~C.5 (conditional on R1--R10) and the finite $V_*$, in the working coordinates
(Remark~3.3 of the record); this is the rider of trap~10, and $M_B$ replaces the record's $M$,
which is not unchanged.  Theorem~\ref{thm:D4} bounds the constructed observable,
$V(h')(0)=J_0V(h')$ with $V(h')$ the limit of Lemma~\ref{aux:notation-wick}(iv); its base term
and its comparison term concern one element of $L^2(\dmu)$, because the two regularizations of a
sharp-time Wick power agree (Lemma~\ref{aux:notation-sharpid} and
Remark~\ref{aux:notation-sharpidrem}(b)).  So $M_B$ and the Euclidean bound need no GAP; through
$9^{p'}$ (Lemma~\ref{aux:notation-hyper}(ii)) and $e^{15\lambda|a_0|}$ they are conditional on
repair R8, and through the comparison term on Lemma~C.4 of the record (repairs R1--R10), through
Lemma~\ref{lem:D2}.  The GAP recorded after Lemma~\ref{aux:notation-wick} enters only where $M_B$ is used
for the record's $V(h)$: at Theorem~5.1, Step~1, and at Corollary~5.2 (\S\ref{sub:cons-51}).
\item It is uniform in $h\in\CE$, hence in $T>0$ and $g\in\CH$ when $h=h_{T,g}$; in $a$,
because the boxes $B_{3\times5}(a)$ (Definition~\ref{aux:notation-ringdef}) all have the shape
$3\times5$ and $S_2$ depends on the shape only (Proposition~\ref{aux:notation-S2}(iii), through
Lemma~C.4, conditional on R1--R10); and in $h'$.
\item The factors:
\begin{itemize}
\item $e^{15\lambda|a_0|}$ is Lemma~C.2 with $|B_{3\times5}|=15$ (conditional on R8);
\item $S_2(B_{3\times5})$ bounds $\|e^{-\Uc_{h^B}}\|_{L^2(\dmu)}$
(Proposition~\ref{aux:notation-S2}(iv));
\item $9^{p'}V_*$ bounds $\|V(h')\|_{L^4(d\varphi_0)}^2$, which equals
$\|V(h')^2\|_{L^2(d\varphi_0)}$, for the constructed $V(h')$ (Lemma~\ref{aux:notation-hyper};
conditional on R8; no GAP);
\item $C_7c_U\overline\Sigma_{3\times5}$ is Lemma~\ref{lem:D2}(a) summed over the rings
(conditional on Lemma~C.4, R1--R10, through Lemma~\ref{lem:D2}),
$\overline\Sigma_{3\times5}$ bounding the sum of $e^{-m_1w_\Delta}$ over the charged squares
outside $B_{3\times5}$;
\item the last factor bounds $\GJS{Q_V}$ as in eq.~(5) of the record \cite[p.~19]{R}, whose
``$=M^2/C_8$'' no longer applies (plan, item n14).
\end{itemize}
\item \emph{Where the transverse case enters.} Every square of rows $-1,0,1$ outside
$B_{3\times5}$ is separated from $R$ only by transverse strips of positive width
(Lemma~\ref{aux:notation-rings}(v)); every side square of the ring $\mathcal L_j$ in a row
$2\le|l|\le j$ admits parallel strips only of width $\le|l|-1<j$, and the bound
$d(\Delta,R)\ge j$ uses a transverse strip for it too
(Remark~\ref{aux:notation-ringrem}(b)). The squares of the rows $|l|=j+1$ of $\mathcal L_j$
are reached by parallel strips of width $j$. Each application with a transverse strip falls
under sentence F-T (\S\ref{sec:purpose}).
\item Theorem~\ref{thm:D4} rests on the printed hypotheses of Theorem~1.1.7; Phase~B does not
remove the weak-coupling restriction (trap~17).
\end{itemize}

\emph{Theorem~5.1 with $M:=M_B$, Corollary~5.2, and \S\S8--14 of the record.}
\begin{itemize}
\item Theorem~5.1 \cite[p.~18]{R}: for every $g\in\CH$ and every real measurable $h'$ with
$\|h'\|_\infty\le1$ and $\supp h'$ in an interval of length $2$ (which lies in some
$[a,a+2]$), $\Om_g\in D(V(h'))$ and $\|V(h')\Om_g\|\le M_B$. Step~5 applies
Theorem~\ref{thm:D4} at $h=h_{T,g}\in\CE$ for every $T>0$; Step~6 lets $T_k\to\infty$ at fixed
$g$. The order ``there is $M_B$ such that for all $g$ and $h'$'' is the content of (M$_2$); the
slab bound implies (M$_2$), and the converse is not claimed (trap~9).  Through
Theorem~\ref{thm:D4}, Theorem~5.1, and with it Corollary~5.2, Lemma~10.2, Theorem~11.1 and
items (1) and (5) of Theorem~12.2, rest on the transverse case of Theorem~1.1.7 at the printed
statement's scope (sentence F-T, \S\ref{sec:purpose}).  They rest on the GAP recorded after
Lemma~\ref{aux:notation-wick} as well, which enters at Step~1 (``$Q_V$ is exactly the
multiplication function $V(h)^2$'', main.tex l.~1278): there the record's $V(h')$ is identified
with the constructed $V(h')$ of Lemma~\ref{aux:notation-wick}(iv), to which
Theorem~\ref{thm:D4} applies (\S\ref{sub:cons-51}).  Lemma~\ref{aux:notation-sharpid} and
Remark~\ref{aux:notation-sharpidrem} identify the two constructions of the note with each other;
they do not reach the record's $V(h)$.
\item Corollary~5.2 \cite[p.~19]{R}: its number $362\,880\,C_8K_1K_2^{16}\lambda^2$ retires
with $C_8$. For $\mathscr P(\xi)=\xi^4$ ($p'=2$, $a_4=1$, the other $a_n=0$) the formula above
gives
\[
 M_B^2=81\,S_2(B_{3\times5})\,V_*+362\,880\,C_7\,c_U\,\overline\Sigma_{3\times5}\,
 K_1K_2^{16}\lambda^2,
\]
with $c_U=24\lambda K_1K_2^8$ and $V_*\le24\lambda^2\iint_{[-1,1]^2}C(x-y)^4\,dx\,dy$.  As a
bound for the record's $V(h)$, Corollary~5.2 inherits from Theorem~5.1 the transverse case (F-T)
and the GAP at Step~1; $M_B$ itself needs no GAP.
\item \S\S8--14 consume the local moment only through Lemma~10.2 \cite[p.~28]{R}, which applies
Theorem~5.1 at the cutoffs $g_s\in\CH$ of an affine path and at $h_j$ supported in
$[j-1,j+1]$; with $M:=M_B$, which is uniform in $g$, Lemma~10.2 holds verbatim (cf.\ trap~11).
The constants $A$, $B'$, $A_f$, $K_f$ and $s_0$ do not leave \S7 (gate GA6,
\S\ref{sec:cons}).
\item In the original coordinates the conclusions are transported by Remark~3.3 of the
record; for the local moment this means supports of length $2/s$, or a partition of unity into
at most $\lceil s\rceil+1$ pieces and the constant $(\lceil s\rceil+1)M_B$, where
$s=m_0/M_*$.
\end{itemize}

\subsection{Boundary cases}\label{sec:quant-boundary}

Squares, rows, columns, $l_s$ and the charged squares $\mathcal D(h,B)$ are those of
\S\S\ref{sec:notation}--\ref{sec:d2}: $\Delta_{l,k}=[l-\frac12,l+\frac12)\times
[k-\frac12,k+\frac12)$, $l_s=\lfloor s+\frac12\rfloor$, and $\mathcal D(h,B)$ the squares
$\Delta\not\subset B$ with $\int_\Delta h>0$.

\emph{(1) $T<1$: a slab thinner than one row.} $h_{T,g}=\mathbf 1_{[-T,T]}\otimes g$ charges at
most the rows $l$ with $|[l-\frac12,l+\frac12)\cap[-T,T]|>0$, i.e.\ $|l|<T+\frac12$: row $0$
alone for $T\le\frac12$, rows $-1,0,1$ for $\frac12<T<1$. In Theorem~\ref{thm:D3} with
$l_s=0$, the box $B$ (rows $-1,0,1$, all columns meeting $\supp g\cup I$ in positive measure)
then contains every charged square, $\mathcal D(h_{T,g},B)=\varnothing$, $h^B=h_{T,g}$, and the
bound is the base term $A\|F\|_{L^2}$ alone. For other $l_s$ the charged squares outside $B$ lie
at row distance $\ge2$ from $l_s$ and the parallel strips of the general case apply.
Theorem~\ref{thm:D3} needs no lower bound on $T$. In Theorem~\ref{thm:D4} with $T\le\frac12$
the only charged row is $0$, and every charged square outside $B_{3\times5}$ lies in row $0$
and is reached only by a transverse strip: the transverse case is used even for the thinnest
slabs. (The fallback Lemma~L of the plan needs $T\ge t_*=1/m_0$; Theorems~\ref{thm:D3}
and~\ref{thm:D4} do not.)

\emph{(2) $T$ not an integer.} The rows at the edge of the slab are charged only in part, so
$h_{T,g}\mathbf 1_\Delta\ne g$ on such squares; $c_U$ uses only
$\|h\mathbf 1_\Delta\|_2\le1$ (Lemma~\ref{aux:notation-cU}(i)), and no constant depends on $T$.
If $T\in\frac12+\Z$, the row $l=T+\frac12$ meets $[-T,T]$ only in the null set $\{x^0=T\}$, is
not charged, and carries $U_\Delta=0$; it needs no strip (Remark~\ref{aux:d2-uses}(3)).

\emph{(3) $s$ on a lattice line.} If $s=l-\frac12$ with $l\in\Z$, then $l_s=l$
(Definition~\ref{aux:notation-squares}): the observable is localized in row $l$, although the
line $x^0=s$ is also the upper edge of row $l-1$. The box $B$ of Theorem~\ref{thm:D3} contains
rows $l-1$, $l$, $l+1$, hence both rows touching the line. A charged square in row $l'$ with
$|l'-l|\ge2$ is separated from the row-$l$ squares by the parallel strip between
$x^0=\min\{l,l'\}+\frac12$ and $x^0=\max\{l,l'\}-\frac12$, of width $|l'-l|-1$; it separates the
closed squares, in particular the line $x^0=s$ itself. For $l'\le l-2$ the line lies on the
boundary of that strip; for $l'\ge l+2$ the true distance is one unit larger than the width
used. The constants are unchanged.

\emph{(4) $N_g=1$, and $N_g=0$.} If $\supp g$ meets one column in positive measure, the charged
squares outside $B$ lie in that column, in rows $l$ with $|l-l_s|\ge2$, at most two at each
strip width $j=|l-l_s|-1\ge1$, and
$B'=2C_7c_Ue^{-m_1}(1-e^{-m_1})^{-1}$. If $g=0$ almost everywhere, no square is charged,
$dq_{h_{T,g}}=\dmu$ (since $\Uc_{h_{T,g}}=0$ and $Z=1$), the comparison term vanishes, and the
theorem is the bound $|\int F\,d\mu_0|\le\|F\|_{L^2}\le A\|F\|_{L^2}$ ($S_2\ge1$). The count
must be taken with the charged squares of Lemma~\ref{lem:D2}(a). With the index set
$\{\Delta\not\subset B:\Delta\cap\supp h\ne\varnothing\}$ of its form (b), a column that meets
$\supp g$ only in a null set would contribute squares in the rows $l_s-1$, $l_s$, $l_s+1$
outside $B$, which only transverse strips separate from $R$; for instance
$g:=\mathbf 1_{[0,1/10]}+\mathbf 1_{\{3\}}\in\CH$ has $N_g=1$ while column $3$ meets
$\supp g=[0,\frac1{10}]\cup\{3\}$. Such columns carry $U_\Delta=0$, and form (a) omits them
(Remark~\ref{aux:d2-uses}(3), which records this example).

\emph{(5) $I$ meeting one column only.} Then $N^\square_I=1$, every cylinder monomial of $F$ has
a single localized piece, all its factors in one square $\Delta_{l_s,k}$, and
$\GJS F$ is bounded by (N.a) of \S\ref{sec:notation} with count $1$. If $\supp g$ meets only the
same column, $N_B=1$, $|B|=3$ and $A=e^{3\lambda|a_0|}S_2(B)$. If $I$ touches a second column in
a null set (e.g.\ $I=[k-\frac12,k+\frac12]$), the pieces there vanish in $L^2$ and give no
localization square. On the record's cells $[j,j+1)$ the count can differ
(Remark~\ref{aux:notation-counts}(ii)); the note counts on the columns.

\emph{(6) $p'=1$.} $\bar n=3$; the pieces $U_{\Delta,n}$ have $n\in\{1,2\}$, the factors of $Q_V$
have degrees $\le2<3$, and its local degrees are $\le4$, which (1.1.7) permits (the bound is on
the individual $n_i$). $c_U=\lambda K_1(|a_1|K_2^2+2|a_2|K_2^4)$, $9^{p'}=9$, and
$9K_1K_2^{8p'}(4p')!=216\,K_1K_2^8$. Nothing degenerates.

\emph{(7) Ring $k=1$.} $\mathcal L_1$ has $8\cdot1+12=20$ squares
(Lemma~\ref{aux:notation-rings}(iii)): the fourteen squares of rows $\pm2$ in the columns
$k_a-2,\dots,k_a+4$, reached by parallel strips of width $1$; and the six squares of rows
$-1,0,1$ in the columns $k_a-2$ and $k_a+4$, reached by transverse strips of width $1$ and by no
parallel strip of positive width. The four corners (rows $\pm2$ in the two side columns) admit
both. The ring contributes the term $20e^{-m_1}$ of $\overline\Sigma_{3\times5}$, of which
$6e^{-m_1}$ rest on the transverse case alone.

\emph{(8) $a$ such that $[a,a+2]$ meets three columns.} $[a,a+2]$ meets exactly the three
columns $k_a,k_a+1,k_a+2$ (Lemma~\ref{aux:notation-rings}(i)); it meets all three in positive
measure unless $a\in\frac12+\Z$, when it meets $k_a+2$ only in the point $a+2$. In that case
the pieces of $Q_V$ in column $k_a+2$ vanish in $L^2$ and $R$ lies in two columns, while
$B_{3\times5}(a)$ is built on all three, which only enlarges the box: the base term and the ring
count are unchanged. Under the ``positive measure'' reading the box would be $3\times4$ for
$a\in\frac12+\Z$; its rings have $(3+2j)(4+2j)-(1+2j)(2+2j)=8j+10\le8j+12$ squares, each at
strip width $\ge j$ from $R$ by Lemma~\ref{aux:notation-strip}, so that
$\overline\Sigma_{3\times5}$ still bounds their sum; its base
term has $e^{12\lambda|a_0|}\le e^{15\lambda|a_0|}$, and its $S_2$ is at most $S_2(B_{3\times5})$
(Proposition~\ref{aux:notation-S2}(ii), the $3\times4$ box lying in a $3\times5$ box). So $M_B$
is valid under both readings, and the record's count in Theorem~5.1, Step~2 (``at most three
half-open unit spatial cells \dots\ in positive measure'') holds on the columns as well.  (The
ring counts $8j+10$ with strip widths $\ge j$ for $j\le12$, and the rows charged by a slab in case
(1), are also checked in exact integer and rational arithmetic by \texttt{parts/check\_d1.py},
parts (4)--(5); \emph{numerical}, run record \texttt{parts/}\allowbreak\texttt{check\_d1.output.txt}, 3 November 2026.)

\emph{(9) The term $n=0$ of $\mathscr P$.} $U_{\Delta,0}=\lambda a_0\int_\Delta h$ is a
constant, with zero covariance under every $dq_{h'}$ (Lemma~\ref{aux:notation-cU}(iii)), and
$c_U$ starts at $n=1$. The coefficient $a_0$ does not change $dq_h$. As a coefficient of the
interaction it enters $A$ and $M_B$ only through $e^{3\lambda|a_0|N_B}$ and $e^{15\lambda|a_0|}$
(Lemma~C.2, conditional on R8) and through $S_2$, an artefact of bounding
$\|e^{-\Uc_{h^B}}\|_2$ and $Z(h^B)$ separately; with $\mathscr P-a_0$ in place of $\mathscr P$
every $dq_h$ is the same and those factors disappear. As a coefficient of the observable it
enters $M_B$ through $V_*$ and $\GJS{Q_V}$, where the term $n=0$ stays: $V(h')$ contains
$\lambda a_0\int h'$, and the sums $\sum_{n=0}^{2p'}$ include it.

\emph{(10) $\lambda=0$.} Then $\Uc_h=0$, $dq_h=\dmu$ for every $h$, both sides of
Lemma~\ref{lem:D2} vanish, $c_U=V_*=0$, $S_2=1$, and $M_B=0=\|V(h')\Om_g\|$.

\emph{(11) Lemma~\ref{lem:D1} at its edges.} For $h_1=h_2$ both sides of
\eqref{eq:d1-identity} vanish ($\Uc_{h_2-h_1}=0$); for constant $F$ the covariance vanishes and
$\ev{F}{h}$ is constant; at $\alpha\in\{0,1\}$ the derivatives are one-sided, and Step~3 of the
proof covers them, the dominator not depending on $\alpha$.
